\begin{recipe}
[% 
    preparationtime = {\unit[45]{min}},
    bakingtime = {\unit[30]{min}},
    bakingtemperature = {\protect\bakingtemperature{
        topbottomheat=\unit[200]{\textcelcius}
    }},
    portion = {\portion{4}},
    %source = {},
    calory = {295kcal (645kcal) | 12g (19g) Protein | 28g (105g) KH | 16g (17g) Fett},
    price = {1,70€ (2,30€)},
    created = {17.05.2026},
    %rating = {1},
]
{Gerösteter Blumenkohl in Kurkuma-Kefir}

\graph{% pictures
    big=pic/hauptgericht/15_geroesteter-blumenkohl-kurkuma-kefir
}

\introduction{%
Gerösteter Blumenkohl mit würziger Kurkuma-Kefir-Sauce und heißem Chili-Gewürzöl. Cremig, säuerlich, leicht scharf und ziemlich gut, wenn man Blumenkohl mal nicht einfach nur traurig dämpfen will.
}

\ingredients{
    900 g & Blumenkohl\index{Gemüse \& Pilze!Blumenkohl}\\
    1 TL & Garam Masala\index{Kräuter \& Gewürze!Garam Masala}\\
    & Meersalz\index{Kräuter \& Gewürze!Salz}\\
    60 ml & Traubenkernöl oder neutrales Öl\index{Öle, Saucen \& Würzmittel!Pflanzenöl}\\
    150 g & Rote Zwiebeln\index{Gemüse \& Pilze!Zwiebeln}\\
    \nicefrac{1}{2} TL & Kurkuma\index{Kräuter \& Gewürze!Kurkuma}\\
    \nicefrac{1}{4} TL & Chilipulver (optional)\index{Kräuter \& Gewürze!Chili}\\
    30 g & Kichererbsenmehl\index{Hülsenfrüchte!Kichererbsenmehl}\\
    480 ml & Kefir oder Buttermilch\index{Milchprodukte \& Alternativen!Buttermilch}\index{Milchprodukte \& Alternativen!Kefir}\\
    \nicefrac{1}{2} TL & Kreuzkümmelsamen\index{Kräuter \& Gewürze!Kreuzkümmel}\\
    \nicefrac{1}{2} TL & Schwarze oder braune Senfkörner\index{Kräuter \& Gewürze!Senfkörner}\\
    1 TL & Rote Chiliflocken\index{Kräuter \& Gewürze!Chili}\\
    2 EL & Koriandergrün oder Petersilie\index{Kräuter \& Gewürze!Koriander}\index{Kräuter \& Gewürze!Petersilie}
}

\preparation{%
    \step%
    Backofen auf 200 °C Ober-/Unterhitze vorheizen.
    
    \step%
    Blumenkohl putzen und in Röschen schneiden. Mit Garam Masala, Salz und 1 EL Öl vermengen und auf einem Backblech oder in einer Auflaufform verteilen.
    
    \step%
    Blumenkohl 20--30 Minuten rösten, bis er goldbraun ist und dunkle Röstaromen bekommt. Nach der Hälfte der Zeit einmal wenden.
    
    \step%
    Währenddessen Zwiebeln schälen und fein hacken. In einem Topf 1 EL Öl erhitzen und die Zwiebeln 4--5 Minuten glasig dünsten.
    
    \step%
    Kurkuma und optional Chilipulver zugeben und etwa 30 Sekunden mitrösten.
    
    \step%
    Hitze reduzieren, Kichererbsenmehl einrühren und 2--3 Minuten unter Rühren garen.
    
    \step%
    Hitze weiter reduzieren und den Kefir langsam unter ständigem Rühren zugeben. Die Sauce nur leicht köcheln lassen, bis sie etwas eindickt.
    
    \step%
    Gerösteten Blumenkohl unterheben, vom Herd nehmen und mit Salz abschmecken.
    
    \step%
    In einem kleinen Topf das restliche Öl stark erhitzen. Kreuzkümmelsamen und Senfkörner 30--45 Sekunden darin rösten, bis sie knistern und anfangen zu bräunen.
    
    \step%
    Vom Herd nehmen, Chiliflocken einrühren und das heiße Gewürzöl über den Blumenkohl geben.
    
    \step%
    Mit Koriander oder Petersilie garnieren und warm servieren.
}

\suggestion[Beilage]{%
    Dazu passt Reis sehr gut. Die Nährwerte und der Preis in Klammern beziehen sich auf das Gericht inklusive 100 g trockenem Reis pro Portion.
}

\hint{
    Kefir oder Buttermilch nicht stark kochen lassen, sonst kann die Sauce gerinnen. Also eher sanft erhitzen und rühren, auch wenn Geduld hier leider wieder nervig wichtig ist.
}

\end{recipe}
